\section{Aula 3}

\begin{question}
    Como se comporta $\Theta(p)$ para $p > p_c$?
\end{question}

\begin{question}
    Quantas componentes infinitas existem em $\Z^d$?
\end{question}

\begin{theorem}
Para todo $p \in [0,1]$ existe no máximo uma componente infinita em $(\Z^d)_p$ q.c.
\end{theorem}

\begin{proof}
    A primeira coisa a notar é que a medida é invariante por translação. Definimos
    $N = \#\{\text{componentes conexas}\}$, como $N$ também é invariante por translação. Por alguma 
    lei 0-1 de ergocidade, segue que $N$ é constante q.c.

    Primeiro a gente prova que $N \in \{0,1,\infty\}$. Seja $N_m$ o número de componentes
    infinitas que tocam o bordo de $B_m$. Note que $N_m$ é independente da configuração da bola $B_m$ em si.
    Então a gente pode brincar de abrir e fechar as arestas dessa bola.
    Como $N_m \to N$ e estamos supondo $N = k \not \in \{0,1,\infty\}$, existe $m_0$ tal que 
    $\Pr(N_{m_0} = k) \geq 1 - \epsilon$. Agora, fixado tal $m_0$ vamos definir $\cE_0$ o evento 
    que todas as arestas de $B_{m_0}$ estão fechados e $\cE_1$ o evento que todas estão abertas.
    Agora segue que 
    $$\Pr(N_m \geq k \cap \cE_1) > 0.$$
    mas se isso acontece, teríamos $\Pr\{N = 1\}>0$ (pois $N \leq k$ q.c).
    Para mim isso já é a prova completa, por algum motivo o Augusto definiu $\cE_0$ e eu não sei exatamente o 
    porquê.

    A segunda parte é bem bonitinha. Para mostrar que $\Pr(N = \infty) = 0$ fazemos uso da bola de raio $m$ mas 
    em $L^1$. O augusto mostra na verdade que $\Pr(N \geq 3) = 0$ e de maneira bem esperta. Assim como antes, 
    se essa probabilidade é positiva, é possivel encontrar uma $L^1$-bola de raio $m$ de que 
    toca em pelo menos três componentes infinitas distintas (pelas bordas). Definimos então um ponto de trifurcação da bola 
    um vértice que se removido, faz aparecer três componentes infinitas distintas. Pela bola que usamos e supondo 
    $\Pr(N \geq 3) > 0$, alterando somente finitas arestas dentro da bola, vemos que com probabilidade positiva, um ponto é de trifurcação.

    Um lema esperto (e brincando com árvores) mostra que podem existir no máximo $|\partial B_m| - 2$ tais pontos. Mas, 
    por linearidade da esperança, em média existem $\alpha|B_m|$ deles. Fazendo $m \to infty$ achamos a bela contradição.
\end{proof}

Depois voltamos a nos interessar pela continuidade de $\Theta(p)$.
\begin{question}
    Será $\Theta(p)$ contínua em algum intervalos? 
\end{question}
\begin{theorem}
    $\Theta(p)$ é crescente, identicamente 0 abaixo de $p_c$, 
    pelo menos contínua a direita em $p_c$ e contínua para $p>p_c$.
\end{theorem}
\begin{proof}
    As primeiras duas são claras, a terceira segue pois $\Theta(p) = \lim_{n \to \infty} \Theta_n(p)$ (lembre que essa é a probabilidade 
    da origem se conectar ao bordo da bola de raio $n$). Como $\Theta_n{p}$ é contínua e crescente e a sequência 
    é decrescente, algum antigo resultade de análise conclui que $\Theta(p)$ é cadlag (nem lembro como prova isso).
    
    A parte mais interessante é a última. Para isso a gente considera o coupling das arestas com variáveis uniformes $[0,1]$. 
    E percebemos que $(\Z^d)_p$ é pegar somente as arestas $e$ tais que $U_e \leq p$. Isso define uma família de percolações na verdade
    e vemos que estamos interessados em $\lim_{p'\to p} \Theta(p) - \Theta(p')$ onde imaginamos $\Theta(p')$ sendo as probabilidades na família $(\Z^d)_{p'}$.
    Agora, se esse limite fosse diferente de $0$, isso se traduziria nesse coupling que 
    \[
        \Pr(0 \leftrightarrow_p \infty \land 0 \not \leftrightarrow_{p'} \infty \forall p'<p) > 0
    \]
    Mas em $p$ existe um cluster infinito $C$ conectando o $0$. Em $p_c < p' < p$ esse cluster infinito já existe também mas não toca o $0$ (denotamos de $C^{p'}$).
    Como o número de arestas que conecta $0$ a $C^{p'}$ em $C$ são finitas e cada uma ocorre aparece antes de $p$ q.c, existe algum $\alpha < p$ tal que $\C^{\alpha} = C$
    e isso é um absurdo, completando a prova.
\end{proof}

